% Propietario conservado: /Users/ruben/Documents/ChatGPT/jueces y controles/REVISION_ARGUMENTAL_APERTURAS_CIERRES_20260915/II/manuscript_es/sections/10c_boltzmann_desarrollo.tex
% Fragmento aditivo. No sustituye 10b_boltzmann.tex.
% Procedencia y distinción entre recuperación y formalización:
% 10c_boltzmann_procedencia.md. Etiquetas locales prefijadas ii-kb-aut.

\subsection{Memoria genealógica, reducción informacional y transducción térmica}
\label{ii:sec:ii-kb-aut-desarrollo}

La construcción térmica utiliza dos publicaciones del mismo desarrollo
APP--TRIT--TPK. La primera conserva los estados, las rutas admisibles y
la información del registro; la segunda proporciona la acción orientada
y la duración del retorno. Su composición determina una energía por
decisión trítica. Las coordenadas de temperatura y energía se introducen
después mediante aplicaciones dimensionales explícitas. Esta secuencia
permite identificar qué información se conserva, qué información publica
un lector y qué magnitud convierte esa información en entropía dimensional.

\subsubsection{Lector, memoria y recuperación del estado}

Fijemos un nivel finito del desarrollo y denotemos por \(\Omega\) el
conjunto de estados enriquecidos admisibles. El registro conserva hoja,
orientación, residuo, cociente, acarreo, ruta, fase y memoria en los
dominios donde interviene cada coordenada. Sean
\[
 q:\Omega\longrightarrow Y,\qquad
 M:\Omega\longrightarrow\mathcal M
\]
la publicación visible y la memoria retenida, respectivamente. Para una
distribución \(p\) sobre \(\Omega\), escribimos
\(\mathsf H(p)=-\sum_xp_x\log p_x\), con \(0\log0=0\).
El logaritmo se aplica a probabilidades adimensionales.

\begin{definition}[Recuperación genealógica de una publicación]
El par \(\widehat q=(q,M)\) es un lector recuperable sobre el soporte
\(\Omega_p\) cuando existe
\[
 d:\widehat q(\Omega_p)\longrightarrow\Omega_p,
 \qquad d\circ\widehat q=\operatorname{id}_{\Omega_p}.
\]
La información de fibra de \(q\) es
\(\mathsf H_{\mathrm{fib}}(q;p)=\mathsf H(X\mid q(X))\), donde
\(X\) tiene ley \(p\).
\end{definition}

\begin{proposition}[Identidad de recuperación y coste de la reducción]
Para un lector recuperable,
\begin{align}
 \mathsf H(X\mid q(X),M(X))&=0,\\
 \mathsf H(X)&=\mathsf H(q(X))+\mathsf H(X\mid q(X)),\\
 \mathsf H(X\mid q(X))&=\mathsf H(M(X)\mid q(X)).
 \label{ii:eq:ii-kb-aut-fibra}
\end{align}
\end{proposition}
\begin{proof}
La aplicación \(d\) recupera \(X\) del par \((q(X),M(X))\),
por lo que la primera entropía condicional es nula. Como \(q\) es
función de \(X\), la regla de la cadena da la segunda igualdad.
Para la tercera, se desarrolla de dos maneras
\(\mathsf H(X,M(X)\mid q(X))\). Una proporciona
\(\mathsf H(X\mid q(X))\), pues \(M\) es función de \(X\);
la otra proporciona
\(\mathsf H(M(X)\mid q(X))+\mathsf H(X\mid q(X),M(X))\).
\end{proof}

La división euclídea de las hojas de APP muestra el mecanismo
aritmético elemental. Con \(i,j\in\{1,\ldots,9\}\), se conserva
\[
 i+j=9q^+_{ij}+R^+_{ij},\qquad
 ij=9q^\times_{ij}+R^\times_{ij},\qquad
 1\leq R^\sigma_{ij}\leq9\quad(\sigma\in\{+,\times\}).
\]
El par residuo--cociente recupera el valor entero de la operación.
Al sumar las diferencias sobre las ochenta y una posiciones se obtiene
\[
 2025-810=54+9\cdot129.
\]
La memoria del cociente restaura la contribución que no publica la
reducción residual. La recuperabilidad del estado completo requiere,
además, conservar las coordenadas de posición y de ruta pertinentes;
el par residuo--cociente identifica aquí el resultado de la operación.

\subsubsection{Actualización reversible, decisión trítica y borrado}

Sea \(U:\Omega\to\Omega\) la actualización TPK en un dominio donde
el transporte enriquecido posee inversa. La ley transportada es
\(p'_y=p_{U^{-1}y}\). La biyección conserva el multiconjunto de
probabilidades y, por tanto,
\begin{equation}
 \mathsf H(U(X))=\mathsf H(X),\qquad
 \mathsf H(X\mid U(X))=0.
 \label{ii:eq:ii-kb-aut-reversible}
\end{equation}
Si la salida observada es \(q(U(X))\), la información no publicada
se cuantifica por \(\mathsf H(X\mid q(U(X)))\). El estado completo,
su publicación y el reinicio de un registro son así tres mapas distintos.

Para los tres estados orientados del TRIT, sea
\(p=(p_{-1},p_0,p_{+1})\) una distribución normalizada. Su información es
\begin{equation}
 I_{\mathrm{TRIT}}(p)=-\sum_{a=-1}^{1}p_a\log p_a,
 \qquad 0\leq I_{\mathrm{TRIT}}(p)\leq\log3.
 \label{ii:eq:ii-kb-aut-trit}
\end{equation}
El extremo superior corresponde a \(p_a=1/3\). En efecto,
\[
 D(p\Vert(1/3,1/3,1/3))
 =\sum_ap_a\log(3p_a)=\log3-I_{\mathrm{TRIT}}(p)\geq0;
\]
la desigualdad se obtiene de la convexidad de \(t\log t\), y la
igualdad caracteriza la distribución uniforme. El extremo inferior
corresponde a una distribución concentrada en un estado.

La realización ideal de Landauer utiliza una memoria con estados lógicos
energéticamente degenerados, un entorno isotermo de temperatura \(\Theta\)
y un balance que conserva la salida del mapa lógico y contabiliza los
registros auxiliares. En ese modelo, la disminución de entropía lógica es
\(k_B[\mathsf H(X)-\mathsf H(F(X))]
=k_B\mathsf H(X\mid F(X))\) para un mapa determinista \(F\).
El balance entrópico exige transferir al entorno al menos esa cantidad;
su expresión energética es \(Q\geq k_B\Theta\mathsf H(X\mid F(X))\).
La distinción entre evolución lógica reversible y borrado físico
corresponde a la realización de Landauer--Bennett \cite{ii:bennett2003}.
La actualización invertible
\eqref{ii:eq:ii-kb-aut-reversible} tiene contribución lógica nula; el
reinicio de una terna uniforme requiere transferir \(\log3\) nats
al entorno y satisface \(Q\geq k_B\Theta\log3\).
La energía de control y la disipación de una realización física
concreta pertenecen a su balance termodinámico adicional.

\subsubsection{Acción, frecuencia y energía de decisión}

Sean \(\mathcal L_{\rm act}\) y \(\mathcal L_T\) las rectas
dimensionales de acción y duración; la recta energética es
\(\mathcal L_E=\mathcal L_{\rm act}\otimes\mathcal L_T^*\).
El desarrollo anterior aporta las secciones positivas
\(\hbar_a\in\mathcal L_{\rm act}\), con
\(a\in\{\mathrm{pre},\mathrm{ret}\}\), y la duración elemental
\(t_0\in\mathcal L_T\). El calendario completo proporciona
\[
 T_{108}=108t_0,\qquad
 \omega_{108}=\frac{2\pi}{T_{108}},\qquad
 h_a=2\pi\hbar_a.
\]
La contracción dimensional entre acción y frecuencia define
\begin{equation}
 \epsilon_a:=\hbar_a\omega_{108}
 =\frac{h_a}{108t_0}
 =\frac{\hbar_a}{t_P},\qquad
 t_P=\frac{54}{\pi}t_0.
 \label{ii:eq:ii-kb-aut-energia}
\end{equation}
El último miembro emplea el mismo período circular reducido que la
sección gravitatoria. Las tres expresiones corresponden a una única
sección de energía, con una orientación \(a\) mantenida en todos
los factores. La información normalizada de una decisión uniforme
produce la sección de energía por nat
\begin{equation}
 \varepsilon_a:=\frac{\epsilon_a}{\log3}\in\mathcal L_E^+.
 \label{ii:eq:ii-kb-aut-energia-nat}
\end{equation}
La acción previamente construida determina la escala de energía al
fijar la duración; el cardinal de alternativas del TRIT determina su
normalización informacional. El valor de \(\log3\) entra por el
conteo y la medida declarados en \eqref{ii:eq:ii-kb-aut-trit}.

\subsubsection{Definición dimensional de Boltzmann y selección de coordenadas}

Sea \(\mathcal L_\Theta\) la recta orientada de temperatura.
El espacio dimensional de entropía es
\(\mathcal L_{\rm ent}=\mathcal L_E\otimes\mathcal L_\Theta^*\).
Una aplicación lineal positiva
\(k_B:\mathcal L_\Theta\to\mathcal L_E\) es, equivalentemente,
un elemento positivo de \(\mathcal L_{\rm ent}\). Por eso el mismo
objeto puede actuar sobre una temperatura o multiplicar una información
adimensional.

\begin{definition}[Normalización térmica de la decisión trítica]
Un par térmico compatible con la sección orientada \(a\) consta de
un transductor positivo \(k_B\) y una sección positiva
\(\Theta_{{\rm clk},a}\in\mathcal L_\Theta\) tales que
\begin{equation}
 k_B(\Theta_{{\rm clk},a})=\varepsilon_a,
 \qquad
 k_B\Theta_{{\rm clk},a}\log3=\epsilon_a.
 \label{ii:eq:ii-kb-aut-par-termico}
\end{equation}
La constante de Boltzmann desempeña la función de transductor
energía--temperatura de la información trítica en esta realización.
\end{definition}

\begin{theorem}[Unicidad relativa a la sección térmica y covariancia dimensional]
Fijada una sección positiva \(\Theta_{{\rm clk},a}\), existe una
única aplicación lineal positiva que satisface
\eqref{ii:eq:ii-kb-aut-par-termico}. Para bases positivas \(u_E,u_\Theta\),
si
\[
 \epsilon_a=\epsilon_a^{\rm num}u_E,\quad
 \Theta_{{\rm clk},a}=\theta_a u_\Theta,\quad
 k_B(u_\Theta)=b\,u_E,
\]
su coordenada es
\begin{equation}
 b=\frac{\epsilon_a^{\rm num}}{\theta_a\log3}.
 \label{ii:eq:ii-kb-aut-coordenada}
\end{equation}
Bajo \(u'_E=s_Eu_E\), \(u'_\Theta=s_\Theta u_\Theta\), con
\(s_E,s_\Theta>0\), se cumple
\[
 b'=\frac{s_\Theta}{s_E}b,\qquad
 \theta'_a=\frac{\theta_a}{s_\Theta},\qquad
 (\epsilon_a^{\rm num})'=\frac{\epsilon_a^{\rm num}}{s_E}.
\]
\end{theorem}
\begin{proof}
Una sección no nula genera una recta. Prescribir su imagen determina
la aplicación lineal única; la positividad de ambas secciones prueba
su positividad. Sustituir las coordenadas en
\eqref{ii:eq:ii-kb-aut-par-termico} da
\(b\theta_a\log3=\epsilon_a^{\rm num}\). Las tres leyes de
transformación se obtienen expresando las mismas secciones y la misma
aplicación en las bases nuevas.
\end{proof}

Una misma constante de Boltzmann debe actuar en las dos orientaciones.
Puesto que \(t_P\) es común, la compatibilidad simultánea equivale a
\begin{equation}
 \frac{\Theta_{{\rm clk},{\rm pre}}}{\Theta_{{\rm clk},{\rm ret}}}
 =\frac{\epsilon_{\rm pre}}{\epsilon_{\rm ret}}
 =\frac{\hbar_{\rm pre}}{\hbar_{\rm ret}}.
 \label{ii:eq:ii-kb-aut-orientaciones}
\end{equation}
En efecto, un isomorfismo de rectas conserva cocientes de secciones.
Recíprocamente, si la igualdad anterior se cumple, el isomorfismo
determinado en una orientación satisface también la normalización de
la otra. Así, el retorno modifica la sección térmica del ciclo en la
misma proporción que la energía, mientras el transductor permanece
común a ambas lecturas.

La ecuación anterior explicita la elección requerida para publicar un
valor individual. Si la sección térmica todavía no ha sido seleccionada
por un lector independiente, para cada \(\lambda>0\) el par
\[
 (k_B,\Theta_{{\rm clk},a})
 \longmapsto(\lambda k_B,\lambda^{-1}\Theta_{{\rm clk},a})
\]
conserva \(\epsilon_a\). Por consiguiente, la presentación de una
coordenada individual internamente seleccionada exige dar el lector de
\(\Theta_{{\rm clk},a}\), las bases y su transporte dimensional.
Las ecuaciones de acción e información permanecen determinadas antes
de esa especificación. La identificación posterior de las bases con
joule y kelvin pertenece a la publicación metrológica del resultado.

\subsubsection{Entropía dimensional e información del estado reducido}

Para una distribución de rutas admitidas \(p\), o para un estado
reducido de matriz de densidad \(\rho\), sean
\[
 s(p)=-\sum_xp_x\log p_x,\qquad
 s(\rho)=-\operatorname{Tr}(\rho\log\rho).
\]
Sus publicaciones dimensionales son
\begin{equation}
 S(p)=k_Bs(p),\qquad S(\rho)=k_Bs(\rho)
 \quad\text{en }\mathcal L_{\rm ent}.
 \label{ii:eq:ii-kb-aut-entropia}
\end{equation}
La evaluación en \(\Theta\) produce una energía
\(\Theta S=k_B\Theta\,s\). Para la terna uniforme a la temperatura
del ciclo, \eqref{ii:eq:ii-kb-aut-par-termico} da exactamente
\[
 S_{\rm TRIT}=k_B\log3,\qquad
 \Theta_{{\rm clk},a}S_{\rm TRIT}=\epsilon_a.
\]
Si \(\rho_A\) es la reducción de la purificación bipartita
construida en la sección areal, \(s(\rho_A)\) mide el entrelazamiento
de esa bipartición. El transporte de unidades conserva sus
autovalores, su orientación y las ocupaciones sectoriales. La escala
areal \(\gamma a_0\ell_P^2\) convierte ocupaciones en área; el
transductor \(k_B\) convierte información en entropía dimensional.
Ambas lecturas utilizan el mismo registro mediante operadores distintos.

\subsubsection{Normalización espectral de las historias térmicas}

Sea \(\mathcal G\) un grafo finito de extensiones admisibles ya
construido, con matriz de adyacencia \(A\). Cada arista \(e\) conserva
sus coordenadas de ruta y recibe un incremento adimensional de acción
\(a_*(e)>0\) del lector correspondiente. Elegida la sección de energía
\(\epsilon_*\), pongamos \(t=\beta\epsilon_*\), donde
\(\beta=(k_B\Theta)^{-1}\). El operador térmico es
\begin{equation}
 (L_tf)(y)=\sum_{e:x\to y}\exp[-t a_*(e)]f(x).
 \label{ii:eq:ii-kb-aut-operador}
\end{equation}
Todos sus exponentes son adimensionales. La condición
\(\rho(L_t)=1\) normaliza el crecimiento ponderado de las historias.

\begin{proposition}[Existencia y unicidad de la presión normalizada]
Supongamos que \(A\) es irreducible, \(\rho(A)>1\), y que
\(0<m\leq a_*(e)\leq M\) para todas las aristas. Existe un único
\(t_*>0\) tal que \(\rho(L_{t_*})=1\), y
\begin{equation}
 \frac{\log\rho(A)}{M}\leq t_*\leq
 \frac{\log\rho(A)}{m}.
 \label{ii:eq:ii-kb-aut-presion}
\end{equation}
\end{proposition}
\begin{proof}
La comparación entrada a entrada proporciona
\(e^{-tM}A\leq L_t\leq e^{-tm}A\) para \(t\geq0\).
La monotonía del radio espectral de matrices no negativas implica
las cotas correspondientes para \(\rho(L_t)\). El radio espectral
es continuo, vale \(\rho(A)>1\) en cero y tiende a cero en infinito.
Para \(s>0\), además,
\(L_{t+s}\leq e^{-sm}L_t\), de modo que
\(\rho(L_{t+s})\leq e^{-sm}\rho(L_t)<\rho(L_t)\).
La continuidad y la disminución estricta prueban existencia y
unicidad. Sustituir el valor uno en las cotas iniciales demuestra
\eqref{ii:eq:ii-kb-aut-presion}.
\end{proof}

La realización local de tres alternativas equiponderadas, representada
por un vértice con tres lazos y \(a_*(e)=1\), satisface
\(\rho(L_t)=3e^{-t}\). Su normalización da \(t_*=\log3\).
Con \(\epsilon_*=\epsilon_a\), resulta
\(\beta_{\rm clk}=\log3/\epsilon_a\), exactamente la relación
\eqref{ii:eq:ii-kb-aut-par-termico}. Este cálculo identifica la
confluencia local trítica. Para un grafo de historias con restricciones
adicionales, la confluencia con el mismo reloj exige comprobar
\[
 \rho\!\left(L_{(\log3/\epsilon_a)\epsilon_*}\right)=1
\]
con sus incrementos efectivos. La normalización local y la presión
del sistema completo conservan así sus respectivos dominios.

La relación obtenida reúne estructura, valor y significado. El registro
determina la información recuperable; las probabilidades sobre las
continuaciones determinan su lectura entrópica; la acción y el período
determinan \(\epsilon_a\); el transductor energía--temperatura
publica la entropía dimensional. El área y el entrelazamiento se
relacionan a través del estado reducido y de sus operadores de borde,
manteniendo visible qué dato produce cada magnitud.
